% basic-test.tex
\documentclass[12pt,a4paper]{article}

\usepackage[utf8]{inputenc}
\usepackage[T1]{fontenc}
\usepackage{amsmath,amssymb}
\usepackage{graphicx}
\usepackage{hyperref}
\usepackage{booktabs}

\title{Basic \LaTeX{} Test Document}
\author{Test Author}
\date{\today}

\begin{document}

\maketitle

\begin{abstract}
This is a basic test document for checking editor support of \LaTeX{} syntax.
It contains sections, lists, math, tables, code, quotes, and references.
\end{abstract}

\tableofcontents
\newpage

\section{Introduction}

This paragraph tests plain text, \textbf{bold}, \textit{italic}, \texttt{monospace},
and inline math such as $E = mc^2$ and $\alpha + \beta = \gamma$.

A second paragraph tests line wrapping and paragraph spacing. The quick brown fox
jumps over the lazy dog. Pack my box with five dozen liquor jugs.

\section{Lists}

\subsection{Itemize}

\begin{itemize}
  \item First item
  \item Second item with \textbf{bold} text
  \item Third item with inline math $x^2 + y^2 = z^2$
\end{itemize}

\subsection{Enumerate}

\begin{enumerate}
  \item First step
  \item Second step
  \item Third step
\end{enumerate}

\section{Mathematics}

\subsection{Inline and Display}

Inline: $a^2 + b^2 = c^2$.

Display:
\[
  \int_{0}^{1} x^2 \, dx = \frac{1}{3}
\]

\subsection{Aligned Equations}

\begin{align}
  f(x) &= x^2 + 2x + 1 \\
       &= (x + 1)^2
\end{align}

\section{Tables}

\begin{table}[h]
  \centering
  \begin{tabular}{lcr}
    \toprule
    Name & Value & Unit \\
    \midrule
    Alpha & 1.0 & m \\
    Beta  & 2.5 & s \\
    Gamma & 3.7 & kg \\
    \bottomrule
  \end{tabular}
  \caption{A simple table.}
  \label{tab:simple}
\end{table}

Table~\ref{tab:simple} shows a basic table.

\section{Code}

\begin{verbatim}
public class Hello {
    public static void main(String[] args) {
        System.out.println("Hello, LaTeX!");
    }
}
\end{verbatim}

\section{Quote}

\begin{quote}
This is a quoted paragraph. It tests indentation and spacing inside a quote environment.
\end{quote}

\section{References}

See Section~\ref{sec:conclusion} for the conclusion.

\section{Conclusion}
\label{sec:conclusion}

This is the end of the basic test document.

\end{document}